\section{String Functions}
\label{sec:cses-string-functions}

\problemstatement{Given a string, compute two classic string functions:
the \textbf{prefix function} (the border/failure array) and the
\textbf{Z-function}. Output both arrays.}

\begin{patternbox}
This problem is a direct exercise in the two linear-time fundamentals.
The \textbf{prefix function} $\pi[i]$ is the longest proper border of the
prefix ending at $i$; the \textbf{Z-function} $z[i]$ is the longest
substring starting at $i$ that matches a prefix of the whole string.
\end{patternbox}

\subsection*{Two views of the same self-similarity}

Both functions measure how a string matches itself. The prefix function
scans left to right, extending or falling back along the failure chain so
each character is charged $O(1)$ amortised. The Z-function maintains the
rightmost matched interval (the ``Z-box'') and seeds each new index from
its mirror inside that box. They are convertible into each other, but
computing each directly with its own linear routine is simplest.

\begin{ideabox}[Prefix Function: Extend or Fall Back]
At index $i$, try to extend the previous border: if $s[i]$ equals
$s[j]$ (where $j=\pi[i-1]$) increment, else follow $j\gets\pi[j-1]$ until a
match or $j=0$. The Z-function mirror-seeds instead. Both amortise to
$O(n)$ because the fallback/expansion total is bounded by $n$.
\end{ideabox}

\subsection*{Algorithm}

\begin{enumerate}
  \item Compute $\pi$: for each $i$, fall back via $\pi$ until $s[i]$
        matches or the border is empty, then set $\pi[i]$.
  \item Compute $z$: maintain $(l,r)$; seed $z[i]$ from $z[i-l]$ if inside
        the box, then extend by direct comparison.
  \item Output both arrays.
\end{enumerate}

\subsection*{Dry run}

\dryrun{$s=$``abacaba''.}

\begin{itemize}
  \item Prefix function: $[0,0,1,0,1,2,3]$ -- e.g.\ $\pi[6]=3$ because
        ``aba'' is the longest border.
  \item Z-function: $[7,0,1,0,3,0,1]$ -- $z[4]=3$ since ``aba'' at index
        $4$ matches the prefix ``aba''.
\end{itemize}

\subsection*{C++ implementation}

\begin{lstlisting}
#include <bits/stdc++.h>
using namespace std;

vector<int> prefixFunction(const string& s) {
    int n = s.size();
    vector<int> pi(n, 0);
    for (int i = 1; i < n; i++) {
        int j = pi[i - 1];
        while (j > 0 && s[i] != s[j]) j = pi[j - 1];
        if (s[i] == s[j]) j++;
        pi[i] = j;
    }
    return pi;
}

vector<int> zFunction(const string& s) {
    int n = s.size();
    vector<int> z(n, 0);
    z[0] = n;
    int l = 0, r = 0;
    for (int i = 1; i < n; i++) {
        if (i < r) z[i] = min(r - i, z[i - l]);
        while (i + z[i] < n && s[z[i]] == s[i + z[i]]) z[i]++;
        if (i + z[i] > r) { l = i; r = i + z[i]; }
    }
    return z;
}

int main() {
    string s;
    cin >> s;
    for (int x : prefixFunction(s)) cout << x << " ";
    cout << "\n";
    for (int x : zFunction(s)) cout << x << " ";
    cout << "\n";
    return 0;
}
\end{lstlisting}

\begin{complexitybox}
\textbf{Time Complexity.} $O(n)$ for each function. \textbf{Space
Complexity.} $O(n)$ per array.
\end{complexitybox}

\begin{takeawaybox}
The prefix function and Z-function are the two workhorses of string
matching; both run in linear time by never re-comparing a character more
than a constant number of times amortised. Every earlier
border/period/matching section is an application of one of them.
\end{takeawaybox}
